\chapter{质点运动学}

\section{速度、加速度、运动学方程和轨道}

\begin{problem}{1.1.1}
一物体做直线运动,它的运动学方程为

\[
x = a t + b t ^ {2} + c t ^ {3}
\]

其中 $a, b, c$ 均为常量. 求：

(1) $t=1\sim2$ 期间的位移,平均速度和平均加速度;

(2) t=2 时的速度和加速度.
\end{problem}
\clearpage


\begin{problem}{1.1.2}
一质点沿 $x$ 方向做直线运动, $t$ 时刻的坐标为 $x = 5t^2 - t^3$ , 式中 $x$ 以米计, $t$ 以秒计. 求:

(1) 第 4 秒内的位移和平均速度；

(2) 第 4 秒内质点所走过的路程.
\end{problem}
\clearpage


\begin{problem}{1.1.3}
一质点 $t = 0$ 时从原点出发，以恒定速率向 $x$ 正方向沿轨道 $x^{2} + (y - r)^{2} = r^{2}$ 运动，其中 $r$ 为常量. $T$ 时又回到原点. 求：

(1) $t=\frac{1}{3}T$ 时的位矢、速度和加速度；

(2) 在 t=0 至 $t=\frac{1}{3}T$ 期间，质点的位移、平均速度和平均加速度.
\end{problem}
\clearpage


\begin{problem}{1.1.4}
一质点从位矢为 $r(0) = 4j$ 的位置以初速度 $\bm{v}(0) = 4i$ 开始运动，其加速度与时间的关系为 $a = 3ti - 2j$ . 所有的长度以米计，时间以秒计. 求：

(1) 经过多长时间质点到达 x 轴；

(2) 到达 x 轴时的位置.
\end{problem}
\clearpage


\begin{problem}{1.1.5}
质点沿直线从静止开始运动,其加速度与时间的关系如图 1.1 所示,作出速度与时间的关系及位置与时间关系图.
\begin{center}
\begin{tikzpicture}[>=Stealth, x=0.72cm, y=1.1cm]
  % Axes
  \draw[->, thick] (0,0) -- (7.4,0) node[below=2pt] {$t/\mathrm{s}$};
  \draw[->, thick] (0,-1.4) -- (0,1.75) node[left=2pt, anchor=north east] {$a/(\mathrm{m\cdot s^{-2}})$};

  % Origin and y-axis ticks
  \node[below left=1pt] at (0,0) {$0$};
  \draw (0,1) -- (-0.08,1) node[left=1pt] {$1$};
  \draw (0,-1) -- (-0.08,-1) node[left=1pt] {$-1$};

  % x-axis ticks and labels
  \foreach \x in {1,3,5} {
    \draw (\x, 0.08) -- (\x, -0.08);
    \node[anchor=north east, inner sep=2pt] at (\x, 0) {$\x$};
  }
  \foreach \x in {2,4,6} {
    \draw (\x, 0.08) -- (\x, -0.08);
    \node[anchor=north west, inner sep=2pt] at (\x, 0) {$\x$};
  }

  % Sawtooth acceleration curve
  \draw[thick] (0,0) -- (1,1);
  \draw[dashed, thick] (1,1) -- (1,-1);
  \draw[thick] (1,-1) -- (3,1);
  \draw[dashed, thick] (3,1) -- (3,-1);
  \draw[thick] (3,-1) -- (5,1);
  \draw[dashed, thick] (5,1) -- (5,-1);
  \draw[thick] (5,-1) -- (6.5,0.5);
\end{tikzpicture}
\end{center}
\end{problem}
\clearpage


\begin{problem}{1.1.6}
一质点以恒定的径向速度分量 $\dot{r} = 4\mathrm{m}\cdot \mathrm{s}^{-1}$ 、横向速度分量与矢径成正比，比例系数为 $3\mathrm{rad / s}$ ，在一平面内运动。当质点离出发点 $4\mathrm{m}$ 时，它的速度及加速度的大小为何值？
\end{problem}
\clearpage


\begin{problem}{1.1.7}
质点的运动学方程为：(1) $r = (4t + 5t^2)i - 5j$ ；(2) $r = (2 + 3t)i - (2t - 3t^2)j$ . 求两运动的轨迹.
\end{problem}
\clearpage


\begin{problem}{1.1.8}
一质点沿直线运动, 其速度为 $v = A t - B x$ , 其中 $A, B$ 均为常数. 当 $t = 0$ 时, $x = 0$ . 试求位置坐标与 $t$ 的函数关系式及 $t = 0$ 时的加速度 $a_0$ .
\end{problem}
\clearpage


\begin{problem}{1.1.9}
某固体燃料火箭的垂直加速度为

\[
a = k \mathrm{e} ^ {- b t} - c v - g
\]

式中 b、c、k 均为常数，v 是获得的垂直速度，g 是重力加速度。在大气中运行时，g 被视为常量。试求火箭发射后 t 秒时的垂直速度。
\end{problem}
\clearpage


\begin{problem}{1.1.10}
图 1.2 中机构包括一根固定的弯成曲线的杆和一根用轴钉可绕 O 点转动的动杆 OC。弯杆曲线的方程为 $r=20\sin2\varphi$ 。木块 A 和 B 钉在一起，A、B 分别沿弯杆和直杆 OC 滑动。若 OC 以恒定角速度 5rad/s 做逆时针转动。求木块 A、B 在 r=10cm 处时的速度和加速度。
\begin{center}
\begin{tikzpicture}[>=Stealth, scale=1]
  % Coordinate axes
  \draw[->, thick] (0,0) -- (4.8,0) node[below] {$x$};
  \draw[->, thick] (0,0) -- (0,4.0) node[left] {$y$};

  % Ground hatching under pivot O
  \draw[thick] (-0.35,-0.22) -- (0.35,-0.22);
  \foreach \x in {-0.28,-0.14,0,0.14,0.28} {
    \draw (\x,-0.22) -- (\x-0.12,-0.36);
  }
  \node[left=4pt] at (0,-0.02) {$O$};

  % Rotation arrow around O
  \draw[->, thick] (125:0.32) arc (125:345:0.32);

  % Curved rod r = 20 sin(2 phi), scaled so 20 cm = 3.6 cm, uniform hollow tube
  \draw[thick, double=white, double distance=3.2pt, samples=90, smooth, domain=6:84]
    plot (\x:{3.6*sin(2*\x)});

  % Straight rod OC at phi = 33 deg
  \def\ang{33}
  \def\rlen{4.55}
  \draw[thick, fill=white]
    (0,0) ++(\ang+90:0.055) -- ++(\ang:\rlen) -- ++(\ang-90:0.11) -- ++(\ang+180:\rlen) -- cycle;

  % Hinge circle at O
  \draw[thick, fill=white] (0,0) circle (0.08);

  % Slider block A-B at intersection r = 3.6 * sin(2*33)
  \pgfmathsetmacro{\rmid}{3.6*sin(2*\ang)}
  \begin{scope}[shift={(\ang:\rmid)}, rotate=\ang]
    \draw[thick, fill=white] (-0.28,-0.16) rectangle (0.28,0.16);
    \node[above left=2pt] at (-0.1,0.16) {$A$};
    \node[below right=2pt] at (0.1,-0.16) {$B$};
  \end{scope}

  % Label C at rod end
  \node[above right=2pt] at (\ang:\rlen) {$C$};

  % Angle phi arc and label in clear open space
  \draw[->] (2.25,0) arc (0:\ang:2.25);
  \node at (24:1.95) {$\varphi$};

  % Dimension line for 20 cm along 45 deg
  \coordinate (Peak) at (45:3.67);
  \draw[thin] (0,0) -- (135:1.95);
  \draw[thin] (Peak) -- ++(135:1.95);
  \draw[<->, thin] (135:1.75) -- node[above left, sloped] {$20\,\mathrm{cm}$} ($(Peak)+(135:1.75)$);
\end{tikzpicture}
\end{center}
\end{problem}
\clearpage


\begin{problem}{1.1.11}
质点以恒定速率 v 沿图示的半径为 R 的圆形轨道运动, 用图 1.3 示的极坐标写出质点在 $0 \leqslant \varphi \leqslant 2\pi$ 的各个位置时的速度和加速度.
\begin{center}
\begin{tikzpicture}[>=Stealth, scale=1]
  \def\R{2.0}
  \def\phiang{28}
  \coordinate (O) at (-\R, 0);
  \coordinate (C) at (0, 0);
  \coordinate (P) at ({2*\phiang}:\R);

  % Circle
  \draw[thick] (C) circle (\R);

  % Polar axis from O through C past right edge
  \draw[thick] (O) -- (\R+0.5, 0);
  \node[left=2pt] at (O) {$O$};

  % Chord r from O to P
  \draw[thick] (O) -- (P) node[midway, above left=1pt] {$r$};

  % Dashed radius from C to P
  \draw[dashed, thick] (C) -- (P);

  % Solid radius arrow R
  \draw[->, thick] (C) -- (-32:\R) node[midway, below left=1pt] {$R$};

  % Tangent velocity vector v at P
  \draw[->, thick] (P) -- ++({2*\phiang+90}:0.85) node[above right=1pt] {$\boldsymbol{v}$};

  % Angle labels
  \node at ($(O)+(\phiang/2:0.75)$) {$\varphi$};
  \node at ($(C)+(0.72*\phiang:0.58)$) {$2\varphi$};
\end{tikzpicture}
\end{center}
\end{problem}
\clearpage


\begin{problem}{1.1.12}
一质点沿心脏线 $r = 2(1 + \cos \varphi)$ 运动，在 $0 < \varphi < 180^{\circ}$ 期间， $\dot{r} = -2$ 。求在此期间质点的速度和加速度。
\end{problem}
\clearpage


\begin{problem}{1.1.13}
对于上题所述的质点运动,求速度平行于 x 轴时的 $\varphi$ 值及在此位置时的速度、加速度的大小.
\end{problem}
\clearpage


\begin{problem}{1.1.14}
离水面高度为 h 的岸上用绳索拉船靠岸, 人以恒定的速率 $v_{0}$ 拉绳, 将船的速度、加速度表示成船离岸的水平距离 L 的函数.
\end{problem}
\clearpage


\begin{problem}{1.1.15}
杆以匀角速 $\omega_{0}$ 绕过其固定端 O 且垂直于杆的轴转动，在 t=0 时，位于 O 点的小珠从相对于杆静止开始沿杆做相对于杆的加速度为 $a_{0}$ 的匀加速运动。求小珠相对于静系的速度、加速度与时间的关系。
\end{problem}
\clearpage


\begin{problem}{1.1.16}
在同一竖直面内的同一水平线上的 A、B 两点分别以 $30^{\circ}$ 、 $60^{\circ}$ 为发射角同时抛出两小球，欲使两小球相遇时都在自己轨道的最高点，已知在 A 点发射的小球发射速率 $v_{AO}=9.8m/s$ ，求 A、B 两点之间的距离 L.
\end{problem}
\clearpage


\begin{problem}{1.1.17}
在小丘上置一靶子，在炮位所在处看靶子的仰角为 $\alpha$ ，炮与靶子间的水平距离为 $L$ ，向目标射击时炮身的仰角为 $\beta$ ，炮弹以什么初速度发射才能击中目标？
\end{problem}
\clearpage


\begin{problem}{1.1.18}
一质点以恒定速率 v=c 沿悬链线

\[
y = \frac {a}{2} \left(\mathrm{e} ^ {\frac {x}{a}} + \mathrm{e} ^ {- \frac {x}{a}}\right)
\]

(其中 a 为正值常数)运动, t=0 时 x=0, 且沿 x 轴负向运动. 求位矢、速度、加速度随时间变化的关系.
\end{problem}
\clearpage


\begin{problem}{1.1.19}
某质点做平面运动, 其径向速度分量 $\dot{r} = c_{1}$ , 横向速度分量 $r\dot{\varphi} = c_{2}, c_{1}, c_{2}$ 均为正值常量. $t = 0$ 时加速度的横向分量 $a_{\varphi} = \frac{c_{1}c_{2}}{b}$ , $b$ 为正值常量, $\varphi = 0$ . 求运动学方程和 $t = 0$ 时加速度的径向分量.
\end{problem}
\clearpage


\begin{problem}{1.1.20}
在 t=0 时,从一点同时以同样大的速率 $v_{0}$ 向各个方向抛出小球.证明:在任一时刻 t,这些小球都位于一个球面上,球的中心以自由落体的加速度下落,球的半径等于 $v_{0}t$ .
\end{problem}
\clearpage


\begin{problem}{1.1.21}
已知一质点做平面运动, 其速率为常量 $c$ , 矢径的角速度的大小为常数 $\omega$ . 求该质点的运动学方程及其轨道. 设 $t=0$ 时, $r=0, \varphi=0$ .
\end{problem}
\clearpage


\begin{problem}{1.1.22}
如图 1.5, 一个动靶以恒定速率 v 在 xy 平面内沿 y=h 的直线飞行, t=0 时在 x=0、y=h 点. 导弹在 t=0 时自原点出发, 以恒定的速率 2v 运动, 速度方向始终指向动靶. 求导弹的运动轨道及导弹击中飞靶的时间.
\begin{center}
\begin{tikzpicture}[>=Stealth, scale=1]
  % Axes
  \draw[->, thick] (0,0) -- (4.8,0) node[above left=2pt] {$x$};
  \draw[->, thick] (0,0) -- (0,3.5) node[left=2pt, anchor=north east] {$y$};

  % Target path y = h and target dot
  \draw[thick] (0,2.7) node[left=2pt] {$h$} -- (4.2,2.7);
  \fill (4.2,2.7) circle (0.08);
  \node[right=3pt] at (4.2,2.7) {$(vt, h)$};

  % Missile position (x,y) and velocity vector
  \coordinate (M) at (2.2, 1.45);
  \coordinate (T) at (4.2, 2.7);
  \draw[dashed, thick] (M) -- (2.2, 2.5);
  \draw[dashed, thick] (M) -- (T);
  \draw[->, thick] (M) -- ($(M)!0.62!(T)$) node[midway, below right=1pt] {$2v$};

  \node[below=2pt] at (M) {$(x,y)$};
  \node at (2.46, 1.98) {$\alpha$};
\end{tikzpicture}
\end{center}
\end{problem}
\clearpage


\begin{problem}{1.1.23}
火车能达到的最大加速度为 $\alpha$ ，最大减速度为 $\beta$ 。求从静止出发到最后静止通过的距离 $a$ 和所需的最短时间。
\end{problem}
\clearpage


\begin{problem}{1.1.24}
直线 LM 在一给定的半径为 R 的圆平面内以等角速 $\omega$ 绕圆周上一点 A 转动. 求此直线与圆的另一交点 B 的速度和加速度的大小见图 1.7.
\begin{center}
\begin{tikzpicture}[>=Stealth, scale=1]
  \def\R{2.0}
  \def\phiang{42}
  \coordinate (O) at (0,0);
  \coordinate (Lend) at (-\R,0);
  \coordinate (Rend) at (\R,0);
  \coordinate (A) at (235:\R);
  \coordinate (B) at (\phiang:\R);

  % Circle
  \draw[thick] (O) circle (\R);

  % Base axis (基轴) and dashed left radius
  \draw[dashed, thick] (Lend) -- (O);
  \draw[thick] (O) -- (\R+0.8, 0) node[below=2pt] {\small 基轴};
  \node[below left=1pt] at (O) {$O$};
  \node[below=2pt] at (0.85, 0) {$R$};

  % Dashed chord from leftmost point to B
  \draw[dashed, thick] (Lend) -- (B);

  % Radius vector OB
  \draw[->, thick] (O) -- (B);
  \node[right=3pt] at (B) {$B$};

  % Angle phi arc
  \draw[->] (1.05,0) arc (0:\phiang:1.05);
  \node at (0.5*\phiang:1.28) {$\varphi$};

  % Line LM at t=0 passing through A and (R,0)
  \draw[thick] ($(A)!-0.45!(Rend)$) -- ($(A)!1.38!(Rend)$);
  \coordinate (Pzero) at ($(A)!0.48!(Rend)$);
  \draw[<-, thick] ($(Pzero)+(-60:0.15)$) -- ++(-60:0.65) node[below right=-2pt] {\small $t=0$ 时的 $LM$};

  % Current line LM passing through A and B
  \draw[thick] ($(A)!-0.32!(B)$) node[below right=1pt] {$L$} -- ($(A)!1.38!(B)$) node[above left=1pt] {$M$};
  \node[below right=2pt] at (A) {$A$};

  % Angles psi
  \node at ($(Lend)+(10:1.12)$) {$\psi$};
  \node at ($(A)+(38:1.25)$) {$\psi$};
\end{tikzpicture}
\end{center}
\end{problem}
\clearpage


\begin{problem}{1.1.25}
质点沿一平面做曲线运动, 其加速度矢量的作用线在任何时刻都通过一固定点 O. 试证明: 质点加速度的大小为 $a = \pm v \frac{dv}{dr}$ , 其中 v 为质点的速率, r 为质点到 O 点的距离, $\frac{dv}{dr} > 0$ 时取正号, $\frac{dv}{dr} < 0$ 时取负号.
\end{problem}
\clearpage


\section{自然坐标、切向加速度和法向加速度}

\begin{problem}{1.2.1}
一质点静止在半径为 $R$ 的半球面上的最高点. 要使质点脱离半球, 一开始就不沿半球滑下, 必须给予质点的最小水平速度多大?
\end{problem}
\clearpage


\begin{problem}{1.2.2}
一质点的运动学方程为

\[
\begin{array}{l} x = R \sin \omega t + \omega R t, \\ y = R \cos \omega t + R, \end{array}
\]

式中 $\omega, R$ 均为常量. 求质点在 y 的最大值、最小值时的切向加速度、法向加速度和轨道的曲率半径.
\end{problem}
\clearpage


\begin{problem}{1.2.3}
一质点沿抛物线 $y = \frac{1}{2} Ax^2$ 运动，设路程从抛物线顶点开始计算，质点运动的路程与时间的关系为 $s = bt^2 + ct, A, b, c$ 均是正值常量，试求在顶点时质点的切向加速度和法向加速度.
\end{problem}
\clearpage


\begin{problem}{1.2.4}
曲柄 $OA = r$ ，绕定轴 $O$ 以匀角速 $\omega$ 转动，连杆 $AB$ 用铰链与曲柄端点 $A$ 连接，并可在具有铰链的滑套 $N$ 内滑动。当 $\varphi = 0$ 时， $A$ 端位于滑套 $N$ 处。已知 $AB = l > 2r$ ，求当 $\varphi = 0$ 时，连杆上 $B$ 点的速度、加速度的大小，切向加速度、法向加速度和轨道的曲率半径（见图1.8）。
\begin{center}
\begin{tikzpicture}[>=Stealth, scale=1]
  \def\r{1.5}
  \def\phiang{56}
  \coordinate (O) at (0,0);
  \coordinate (A) at (\phiang:\r);
  \coordinate (N) at (\r,0);
  % Direction from A to N
  \pgfmathsetmacro{\dx}{\r - \r*cos(\phiang)}
  \pgfmathsetmacro{\dy}{0 - \r*sin(\phiang)}
  \pgfmathsetmacro{\rodang}{atan2(\dy,\dx)}
  \def\Llen{3.45}
  \coordinate (B) at ($(A)+(\rodang:\Llen)$);

  % Dashed circle of radius r
  \draw[dashed, thick] (O) circle (\r);
  \draw[dashed, thick] (O) -- (N);

  % Ground support under O
  \draw[thick] (-0.28,-0.32) -- (0.28,-0.32);
  \foreach \x in {-0.2,-0.1,0,0.1,0.2} {
    \draw (\x,-0.32) -- (\x-0.1,-0.44);
  }
  \draw[thick] (-0.16,-0.32) -- (O) -- (0.16,-0.32);
  \node[left=4pt] at (O) {$O$};

  % Crank OA (double-line bar)
  \draw[thick, fill=white]
    (O) ++(\phiang+90:0.06) -- ($(A)+(\phiang+90:0.06)$) -- ($(A)+(\phiang-90:0.06)$) -- ($(O)+(\phiang-90:0.06)$) -- cycle;
  \node[above left=2pt] at ($0.42*(O)+0.42*(A)$) {$r$};

  % Connecting rod AB (double-line bar)
  \draw[thick, fill=white]
    (A) ++(\rodang+90:0.045) -- ($(B)+(\rodang+90:0.045)$) -- ($(B)+(\rodang-90:0.045)$) -- ($(A)+(\rodang-90:0.045)$) -- cycle;

  % Hinge pins at O and A
  \draw[thick, fill=white] (O) circle (0.09);
  \draw[thick, fill=white] (A) circle (0.09);
  \node[above=3pt] at (A) {$A$};
  \node[below right=1pt] at (B) {$B$};

  % Slider sleeve N
  \begin{scope}[shift={(N)}, rotate=\rodang]
    \draw[thick, fill=white] (-0.32,-0.12) rectangle (0.32,0.12);
    \draw[thick] (0,0) circle (0.07);
  \end{scope}
  \node[above right=2pt] at ($(N)+(0.18,0.06)$) {$N$};

  % Angle phi and angular velocity omega
  \draw[->] (0.68,0) arc (0:\phiang-3:0.68);
  \node at (0.5*\phiang:0.92) {$\varphi$};
  \draw[->, thick] (\phiang-18:1.05) arc (\phiang-18:\phiang+26:1.05);
  \node at (\phiang+34:1.14) {$\omega$};

  % Dimension L
  \draw[thin] (A) -- ++(\rodang+90:0.95);
  \draw[thin] (B) -- ++(\rodang+90:0.95);
  \draw[<->, thin] ($(A)+(\rodang+90:0.82)$) -- node[fill=white, inner sep=1.5pt] {$L$} ($(B)+(\rodang+90:0.82)$);
\end{tikzpicture}
\end{center}
\end{problem}
\clearpage


\begin{problem}{1.2.5}
质点 $P$ 沿一平面曲线运动, 其加速度矢量的作用线与曲率圆交于 $A$ 点, $PA$ 间距离为 $l$ . 试证: 质点的加速度大小为 $a = \frac{2v^2}{l}$ , 其中 $v$ 是质点的速率.
\end{problem}
\clearpage


\begin{problem}{1.2.6}
某个在平面上运动的质点, 其速度的 $y$ 分量 $\dot{y} = c$ 为常量. 试证其加速度值 $a = \frac{v^3}{c\rho}$ , 式中 $v$ 为速率, $\rho$ 为曲率半径.
\end{problem}
\clearpage


\begin{problem}{1.2.7}
一质点的运动轨道为对数螺旋线 $r = b\mathrm{e}^{k\varphi}$ , 其中 $b, k$ 均为正值常量, $\dot{r} = c$ ( $c$ 为正值常量), $t = 0$ 时, 位于 $r = b, \varphi = 0$ 处. 求质点的速度和加速度的大小以及曲率半径随时间的变化规律.
\end{problem}
\clearpage


\begin{problem}{1.2.8}
质点做平面运动,轨道方程为 $y=y(x)$ ,试证:

\[
\rho = \frac {(1 + y ^ {\prime 2}) ^ {3 / 2}}{| y ^ {\prime \prime} |}
\]

式中 $\rho$ 为曲率半径， $y' = \frac{\mathrm{d}y}{\mathrm{d}x}, y'' = \frac{\mathrm{d}^2y}{\mathrm{d}x^2}$ .
\end{problem}
\clearpage


\begin{problem}{1.2.9}
已知质点的运动轨道为 $r = r(\varphi), (r, \varphi$ 为极坐标)，求轨道上各点的曲率半径.
\end{problem}
\clearpage


\begin{problem}{1.2.10}
求 1.1.10 题所述轨道在 r=10cm 处的曲率半径.
\end{problem}
\clearpage


\section{质点的相对运动}

\begin{problem}{1.3.1}
一直升飞机在离地面 $4.9\mathrm{m}$ 的高度以 $4.9\mathrm{m / s}$ 的恒定水平速度飞行，一包裹从直升飞机上沿水平方向丢出时相对于直升飞机的初速度为 $12\mathrm{m / s}$ ，方向与飞行方向相反。求：

(1) 当包裹落到地面时, 它与直升飞机的水平距离;

(2) 包裹刚要着地时其速度矢量与地面的夹角.
\end{problem}
\clearpage


\begin{problem}{1.3.2}
一列火车在刮着北风的雨中以 $30\mathrm{m / s}$ 的速率(相对于地面)向正南方行驶。一在地面上静止的观察者测得雨滴的路径与竖直线成 $25^{\circ}$ 角，火车中的观察者看到雨水在车窗玻璃上的径迹是竖直向下的。求雨滴相对于地面的速率。
\end{problem}
\clearpage


\begin{problem}{1.3.3}
一架飞机在静止空气中的速率为 60m/s，风的速率为 30m/s，飞机仍能向北以 60m/s 的速率飞行。求：(1) 风向；(2) 机头的指向。
\end{problem}
\clearpage


\begin{problem}{1.3.4}
一飞机带的燃料, 在无风的情况下, 以 $v$ 的速率飞行, 可飞行来回的总路程为 $R$ . 今在有风的情况下, 要向北偏东 $\varphi^{\circ}$ 角的方向去执行任务, 并飞回基地. 若风向为北偏东 $\alpha^{\circ}$ , 风速为 $\bm{u}$ . 问该机能飞离基地多大距离?
\end{problem}
\clearpage


\begin{problem}{1.3.5}
蒸汽轮船以 $v_{1}$ 的速度向南航行, 船上的人看到烟囱中冒出的烟向东方向, 若船速增至 $v_{2}$ , 船上的人看到烟向东北, 若烟从烟囱冒出后立即按风速运动, 求风速.
\end{problem}
\clearpage


\begin{problem}{1.3.6}
一个小环 $M$ 同时套在两根细棒上，棒 $AB$ 以速度 $\bm{v}_{1}$ 沿垂直于 $AB$ 的方向移动，棒 $CD$ 以速度 $\bm{v}_{2}$ 沿垂直于 $CD$ 的方向移动，两棒始终处在同一平面上，两棒夹角为 $\alpha$ ，如图1.13所示.求小环 $M$ 的速度.
\begin{center}
\begin{tikzpicture}[>=Stealth, scale=1]
  \def\ang{42}
  % Rod AB
  \draw[thick] (-2.6,0) node[below=2pt] {$A$} -- (2.4,0) node[below=2pt] {$B$};

  % Rod CD
  \draw[thick] (180+\ang:2.35) node[below right=1pt] {$C$} -- (\ang:2.35) node[right=2pt] {$D$};

  % Ring M at intersection (0,0)
  \draw[thick] (0,0) circle (0.11);
  \node[above left=2pt] at (0,0.05) {$M$};

  % Angle alpha
  \node at (0.5*\ang:0.55) {$\alpha$};

  % Velocity vectors
  \draw[->, thick] (-1.35,0) -- (-1.35,1.05) node[right=2pt] {$\boldsymbol{v}_1$};
  \coordinate (Q) at (180+\ang:1.25);
  \draw[->, thick] (Q) -- ++(\ang-90:1.05) node[right=2pt] {$\boldsymbol{v}_2$};
\end{tikzpicture}
\end{center}
\end{problem}
\clearpage


\begin{problem}{1.3.7}
圆心为 A、B，半径均为 R=5cm 的两个大圆环处在同一平面上。B 环固定，A

环沿着 $AB$ 连线向 $B$ 环运动. 另有一小环 $M$ 同时套在两个大圆环上. 当 $A$ 环运动到 $\alpha = 30^{\circ}$ 时, $A$ 点的速度 $v_{A} = 5\mathrm{cm}\cdot \mathrm{s}^{-1}$ , 加速度 $a_{A} = 0$ . 求此时小环 $M$ 的速度和加速度的大小.
\end{problem}
\clearpage


\begin{problem}{1.3.8}
有一与 $x$ 轴夹角为 $45^{\circ}$ 的直线，以速率 $u$ 垂直于直线的方向向 $x$ 轴负向、 $y$ 轴正向运动。设在 $t = 0$ 时，它与悬链线

\[
y = \frac {a}{2} \left(\mathrm{e} ^ {\frac {x}{a}} + \mathrm{e} ^ {- \frac {x}{a}}\right)
\]

(其中 a 为正值常数)交于 x=0 处, 求此交点的速度作为时间的函数. 另一个交点移动的速度又如何随时间变化?
\end{problem}
\clearpage


\begin{problem}{1.3.9}
用相对运动理论求 1.1.24 题中 B 点相对于直线 LM 的速度和加速度.
\end{problem}
\clearpage


\begin{problem}{1.3.10}
一个半径为 $R$ 的转盘绕通过盘心的垂直盘面的轴以恒定角速度 $\omega$ 转动。两个小孩 $A$ 和 $B$ 相对于圆盘静止于盘的边缘， $\angle AOB = 90^{\circ}$ ，用图画出在图1.18所示位置， $A$ 看到的 $B$ 的速度和加速度，标明大小和方向，分别选用以 $A$ 为原点的平动参考系和转动参考系。
\begin{center}
\begin{tikzpicture}[>=Stealth, scale=1]
  \def\R{1.85}
  \coordinate (O) at (0,0);
  \coordinate (A) at (0,\R);
  \coordinate (B) at (-\R,0);

  % Disk
  \draw[thick] (O) circle (\R);

  % Radii OA and OB
  \draw[thick] (O) -- (A) node[above=2pt] {$A$};
  \draw[thick] (O) -- (B) node[left=2pt] {$B$};
  \node[below right=1pt] at (O) {$O$};

  % Counterclockwise angular velocity arrow omega
  \draw[->, thick] (-45:0.62) arc (-45:155:0.62);
  \node at (45:0.86) {$\omega$};
\end{tikzpicture}
\end{center}
\end{problem}
\clearpage


\begin{problem}{1.3.12}
平面 I 可绕通过固定点 $O_{1}$ (在平面 I 上) 的垂直轴转动, 平面 II 可绕通过 $O_{2}$ (固连于平面 I) 的垂直于两平面的轴转动. 若 $\omega_{1} = \omega_{2}$ , 如图 1.23 所示. 证明在静参考系中平面 II 作平动.

提示: 只要证明平面 Ⅱ 上任一点 A 的绝对速度与 A 的位置无关即可.
\begin{center}
\begin{tikzpicture}[>=Stealth, scale=1]
  % Plane I
  \draw[thick] (0,0) -- (4.3,0) -- (5.3,2.2) -- (1.0,2.2) -- cycle;
  \node at (4.75, 1.9) {$\mathrm{I}$};

  % Plane II (smaller parallelogram on top of I)
  \draw[thick, fill=white] (1.95,0.72) -- (3.55,0.72) -- (3.95,1.5) -- (2.35,1.5) -- cycle;
  \node at (2.28, 0.92) {$\mathrm{II}$};

  % Axis O1
  \draw[thick] (1.3, 1.25) -- (1.3, 2.95);
  \draw[dashed, thick] (1.3, 0) -- (1.3, 1.05);
  \draw[thick] (1.3, -0.55) -- (1.3, 0);
  \node[left=3pt] at (1.3, 1.15) {$O_1$};

  % Rotation arrow omega1 (counterclockwise viewed from above, centered on axis x=1.3)
  \draw[->, thick] ([shift=(145:0.34 and 0.13)]1.3, 2.54) arc (145:395:0.34 and 0.13);
  \node[right=2pt] at (1.68, 2.54) {$\omega_1$};

  % Axis O2
  \draw[dashed, thick] (2.85, 0) -- (2.85, 1.08);
  \draw[thick] (2.85, -0.78) -- (2.85, 0);
  \node[right=3pt] at (2.85, 1.08) {$O_2$};

  % Rotation arrow omega2 (opposite sense: clockwise viewed from above, centered on axis x=2.85)
  \draw[->, thick] ([shift=(35:0.34 and 0.13)]2.85, -0.38) arc (35:-215:0.34 and 0.13);
  \node[right=2pt] at (3.22, -0.40) {$\omega_2$};
\end{tikzpicture}
\end{center}
\end{problem}
\clearpage


\begin{problem}{1.3.13}
一个质点在 $xy$ 平面内的运动学方程为 $x = f(t), y = g(t)$ ，此平面以等角速 $\omega$ 绕过原点在 $xz$ 平面上与 $x, z$ 轴的夹角均为 $45^{\circ}$ 的固定轴转动。求质点对静参考系的速度和加速度（用动坐标系表达）。
\end{problem}
\clearpage


\begin{problem}{1.3.14}
一个质点在 $xy$ 平面内的运动学方程为 $r = \mathrm{e}^{ct}, \varphi = bt(c, b$ 为常量), $r, \varphi$ 为极坐标. 此平面以等角速 $\omega$ 绕固定的 $x$ 轴转动. 求质点的绝对速度和绝对加速度 (见图1.25).
\begin{center}
\begin{tikzpicture}[>=Stealth, scale=1]
  \def\ang{30}
  \def\rlen{3.2}
  \coordinate (O) at (0,0);
  \coordinate (P) at (\ang:\rlen);

  % Axes
  \draw[->, thick] (O) -- (4.2,0) node[right=2pt] {$x$};
  \draw[->, thick] (O) -- (0,3.2) node[left=2pt] {$y$};
  \node[below left=1pt] at (O) {$O$};

  % Position vector line r
  \draw[thick] (O) -- (P) node[midway, above left=1pt] {$\boldsymbol{r}$};
  \node at (0.5*\ang:1.05) {$\varphi$};

  % Unit vectors at P
  \draw[->, thick] (P) -- ++(0.9, 0) node[right=1pt] {$\boldsymbol{i}$};
  \draw[->, thick] (P) -- ++(0, 0.9) node[above=1pt] {$\boldsymbol{j}$};
  \draw[->, thick] (P) -- ++(\ang:0.95) node[above right=0pt] {$\boldsymbol{e}_r$};
  \draw[->, thick] (P) -- ++(\ang+90:0.95) node[above left=0pt] {$\boldsymbol{e}_\varphi$};

  % Right-angle mark between OP and e_phi
  \draw[thin] ($(P)+(\ang+180:0.18)$) -- ++(\ang+90:0.18) -- ++(\ang:0.18);
\end{tikzpicture}
\end{center}
\end{problem}
\clearpage


\begin{problem}{1.3.15}
图 1.26 中机构在图示位置时, 杆 OA 绕 O 点的角速度为 2rad/s, 求此时 C 的角速度 $\omega$ 及 B 点的速度 $v_{B}$ .
\begin{center}
\begin{tikzpicture}[>=Stealth, scale=0.18, every node/.style={font=\small}]
  % Coordinates (1 unit = 1 cm in problem, scaled by 0.18 -> ~5.8 cm wide)
  \coordinate (O) at (0, 14);
  \coordinate (A) at (8, 14);
  \coordinate (D) at (10, 0);
  \coordinate (B) at (24, 2);

  % Wall at O
  \draw[thick] (-1.5, 10.5) -- (-1.5, 17.5);
  \fill[pattern=north east lines] (-2.7, 10.5) rectangle (-1.5, 17.5);
  % Bearing bracket at O
  \draw[thick] (-1.5, 15.6) -- (0, 15.6) arc (90:-90:1.6) -- (-1.5, 12.4);
  \draw[thick, fill=white] (O) circle (0.7);

  % Crank OA (double-line bar)
  \draw[thick] (0, 14.5) -- (8, 14.5);
  \draw[thick] (0, 13.5) -- (8, 13.5);
  \draw[thick, fill=white] (O) circle (0.55);
  \draw[thick, fill=white] (A) circle (0.55);

  % Angular velocity arrow around OA
  \draw[->, thick] (1.6, 16.2) arc (65:-75:2.5);
  \node at (4.4, 12.1) {$2\,\mathrm{rad/s}$};

  % Ground and bearing bracket at D
  \draw[thick] (6.5, -1.6) -- (13.8, -1.6);
  \fill[pattern=north east lines] (6.5, -2.7) rectangle (13.8, -1.6);
  \draw[thick] (8.1, -1.6) -- (8.1, 0) arc (180:0:1.9) -- (11.9, -1.6);
  \draw[thick, fill=white] (D) circle (0.6);

  % Rod AB and sliding sleeve C
  % Direction of AB: angle -36.87 deg (cos=0.8, sin=-0.6)
  % Sleeve spans from s=5.5 to s=14.5 along AB
  \coordinate (S1) at ($(A)!0.275!(B)$);
  \coordinate (S2) at ($(A)!0.725!(B)$);
  \coordinate (perp) at (0.6, 0.8); % unit normal to AB pointing up-right

  % Rod AB
  \draw[thick] ($(A)+0.4*(perp)$) -- ($(S1)+0.4*(perp)$);
  \draw[thick] ($(A)-0.4*(perp)$) -- ($(S1)-0.4*(perp)$);
  \draw[thick] ($(S2)+0.4*(perp)$) -- ($(B)+0.4*(perp)$);
  \draw[thick] ($(S2)-0.4*(perp)$) -- ($(B)-0.4*(perp)$);
  \draw[thick] ($(B)+0.4*(perp)$) -- ($(B)-0.4*(perp)$);

  % Sleeve around AB
  \draw[thick] ($(S1)+0.85*(perp)$) -- ($(S2)+0.85*(perp)$)
            -- ($(S2)-0.85*(perp)$) -- ($(S1)-0.85*(perp)$) -- cycle;
  \draw[dashed] ($(S1)+0.4*(perp)$) -- ($(S2)+0.4*(perp)$);
  \draw[dashed] ($(S1)-0.4*(perp)$) -- ($(S2)-0.4*(perp)$);
  \draw[thick, fill=white] (A) circle (0.55);

  % Rocker C arms from D to sleeve
  \coordinate (M1) at ($(S1)!0.22!(S2) - 0.85*(perp)$);
  \coordinate (M2) at ($(S1)!0.78!(S2) - 0.85*(perp)$);
  \draw[thick] (D) -- (M1);
  \draw[thick] (D) -- (M2);
  \draw[thick, fill=white] (D) circle (0.55);

  % Angular velocity arrow around D
  \draw[->, thick] ($(D)+(-15:2.9)$) arc (-15:125:2.9);
  \node at (10.7, 3.6) {$C$};
  \node at ($(D)+(-12:4.1)$) {$\omega$};

  % Dashed geometric lines
  \draw[dashed] (A) -- (13.5, 14); % horizontal right from A
  \draw[dashed] (A) -- (8, 5.5);   % vertical down from A
  \draw[dashed] (A) -- (D);        % AD
  \draw[dashed] (D) -- (B);        % DB

  % Angle arcs alpha and beta
  \draw (8, 8.5) arc (-90:-81.87:5.5);
  \draw[->] (5.6, 8.3) to[bend right=25] (8.4, 8.8);
  \node[left] at (5.6, 8.5) {$\alpha$};
  \draw (11.5, 14) arc (0:-36.87:3.5);
  \node at (12.6, 12.6) {$\beta$};

  % Dimension: 8 cm
  \draw (0, 16) -- (0, 19.5);
  \draw (8, 15.5) -- (8, 19.5);
  \draw[<->] (0, 18.5) -- node[above] {$8\,\mathrm{cm}$} (8, 18.5);

  % Dimension: 14 cm
  \draw (-4.5, 14) -- (-0.8, 14);
  \draw (-4.5, 0) -- (7.5, 0);
  \draw[<->] (-3.5, 0) -- node[left, rotate=90, anchor=south] {$14\,\mathrm{cm}$} (-3.5, 14);

  % Dimension: 10 cm (horizontal O to D)
  \draw (0, 11.5) -- (0, -5.6);
  \draw (10, -2.9) -- (10, -5.6);
  \draw[<->] (0, -4.8) -- node[below] {$10\,\mathrm{cm}$} (10, -4.8);

  % Dimension: 20 cm (parallel to AB)
  \draw ($(A)+1.2*(perp)$) -- ($(A)+5.8*(perp)$);
  \draw ($(B)+1.2*(perp)$) -- ($(B)+5.8*(perp)$);
  \draw[<->] ($(A)+4.8*(perp)$) -- node[above right, pos=0.5] {$20\,\mathrm{cm}$} ($(B)+4.8*(perp)$);

  % Dimension: 10 cm (perpendicular from D to line AB)
  \coordinate (D_ext) at ($(D) + 18*(0.8, -0.6)$);
  \coordinate (B_ext) at ($(B) + 4.5*(0.8, -0.6)$);
  \draw ($(D)+2.5*(0.8,-0.6)$) -- (D_ext);
  \draw (B) -- (B_ext);
  \draw[<->] ($(D) + 16.5*(0.8, -0.6)$) -- node[below right] {$10\,\mathrm{cm}$} ($(B) + 2.5*(0.8, -0.6)$);

  % Point labels
  \node[above right] at (8.2, 14.4) {$A$};
  \node[above right] at (B) {$B$};
  \draw[->] (1.4, 11.1) -- ($(O)+(0.2,-0.3)$);
  \node[below] at (1.5, 11.1) {$O$};
  \draw[->] (13.0, -3.6) -- ($(D)+(0.3,-0.4)$);
  \node[below right] at (12.8, -3.4) {$D$};
\end{tikzpicture}
\end{center}
\end{problem}
\clearpage


\begin{problem}{1.3.16}
若上题除所给条件以外再加上在图示位置 $OA$ 杆绕 $O$ 点的角加速度为零，求此时 $C$ 的角加速度 $\dot{\omega}$ 及 $B$ 点的加速度 $a_{B}$ .
\begin{center}
\begin{tikzpicture}[>=Stealth, every node/.style={font=\small}]
  \coordinate (O) at (0, 0);
  \fill (O) circle (1.5pt);

  % Dashed reference lines
  \draw[dashed] (O) -- (1.9, 0);
  \draw[dashed] (O) -- (-37:2.2);
  \draw[dashed] (O) -- (0, -2.6);

  % Angle labels
  \draw (0.85, 0) arc (0:-37:0.85);
  \node at (-18:1.15) {$\beta$};
  \node at (-81:1.35) {$\alpha$};

  % Acceleration vectors
  \draw[->, thick] (O) -- (-2.1, 0) node[below left] {$\boldsymbol{a}_{A}$};
  \draw[->, thick] (O) -- (143:1.8) node[above left] {$\boldsymbol{a}_{Ar}$};
  \draw[->, thick] (O) -- (53:1.45) node[above right] {$\boldsymbol{a}_{Ac}$};
  \draw[->, thick] (O) -- (-164:1.65) node[below] {$\boldsymbol{a}_{Aec}$};
  \draw[->, thick] (O) -- (-72:2.05) node[right] {$\boldsymbol{a}_{Aen}$};
\end{tikzpicture}
\end{center}
\end{problem}
\clearpage

